\begin{examplebox}

	\textbf{\textcolor{CExample}{例 1（基础）}}

	\textbf{\textcolor{CExample}{题目：}}
	原命题：若 $x>0$，则 $x+1>0$。写出其否定和否命题，并判断真假。

	\textbf{\textcolor{CExample}{解题步骤：}}
	\begin{description}[style=nextline, leftmargin=2.8em, labelsep=0.5em, itemsep=0.45em]
		\item[\textbf{第一步（识别结构）：}] 确定条件和结论。原命题的条件 $p$：$x>0$；结论 $q$：$x+1>0$。原命题属于全称命题（省略 $\forall\,x\in\mathbb{R}$）。
			\par\textbf{理由：}
			分清 $p$ 和 $q$ 才能正确写出否定和否命题。

		\item[\textbf{第二步（写出否定）：}] 命题的否定：存在 $x>0$ 使得 $x+1\leq 0$，记作
			\[
				\exists x>0,\; x+1\leq 0.
			\]
			逻辑公式：$\lnot(p\to q)\equiv p\land\lnot q$。
			\par\textbf{理由：}
			否定原命题就是找到使条件成立而结论不成立的反例。

		\item[\textbf{第三步（否命题判真假）：}] 否命题：若 $x\leq 0$，则 $x+1\leq 0$，即
			\[
				x\leq 0 \to x+1\leq 0.
			\]
			判断真假：命题的否定为假（因为对任意 $x>0$ 都有 $x+1>0$）；否命题为假（反例 $x=0$，满足 $x\leq0$ 但 $x+1=1>0$）。
			\par\textbf{理由：}
			命题的否定为假因为原命题真；否命题假可举出反例 $x=0$。
	\end{description}

	\textbf{\textcolor{CExample}{关键结论：}}
	\textbf{命题的否定与原命题真假相反；否命题只是改写条件结论，其真假与原命题未必相同。}

	\medskip
	\textbf{\textcolor{CExample}{例 2（稍变形）}}

	\textbf{\textcolor{CExample}{题目：}}
	原命题：若 $a>b$，则 $a^2>b^2$。写出其否定和否命题，并判断真假。

	\textbf{\textcolor{CExample}{解题步骤：}}
	\begin{description}[style=nextline, leftmargin=2.8em, labelsep=0.5em, itemsep=0.45em]
		\item[\textbf{第一步（识别结构）：}] 条件 $p$：$a>b$；结论 $q$：$a^2>b^2$。原命题假。
			\par\textbf{理由：}
			因为 $a>b$ 未必推出 $a^2>b^2$，例如 $a=1,\,b=-2$。

		\item[\textbf{第二步（写出否定）：}] 命题的否定：存在 $a>b$ 但 $a^2\leq b^2$，即
			\[
				\exists a>b,\; a^2\leq b^2.
			\]
			该命题为真（上述 $a=1,\,b=-2$ 即为反例）。
			\par\textbf{理由：}
			否定只需举出一个符合条件的实例即可。

		\item[\textbf{第三步（否命题判真假）：}] 否命题：若 $a\leq b$，则 $a^2\leq b^2$，即
			\[
				a\leq b \to a^2\leq b^2.
			\]
			该否命题为假（反例 $a=-1,\,b=0$，满足 $a\leq b$ 但 $a^2=1>0$）。
			\par\textbf{理由：}
			否命题的真假与原命题无关，需独立判断。
	\end{description}

	\textbf{\textcolor{CExample}{关键结论：}}
	\textbf{原命题假时其否定真，但否命题可能真也可能假；本例中原命题假，否命题也假，说明否命题与原命题无必然真值关系。}

	\medskip
	\textbf{\textcolor{CExample}{例 3（综合应用）}}

	\textbf{\textcolor{CExample}{题目：}}
	已知命题 $p$：若整数 $n$ 的平方是偶数，则 $n$ 是偶数。写出其逆否命题，并用等价性证明 $p$ 为真。

	\textbf{\textcolor{CExample}{解题步骤：}}
	\begin{description}[style=nextline, leftmargin=2.8em, labelsep=0.5em, itemsep=0.45em]
		\item[\textbf{第一步（写出逆否命题）：}] 原命题 $p$ 的逆否命题为 $\lnot q\to\lnot p$（若 $n$ 不是偶数，则 $n^2$ 不是偶数），即“若 $n$ 为奇数，则 $n^2$ 为奇数”。
			\par\textbf{理由：}
			逆否命题是将原命题的条件与结论交换并分别取否定。

		\item[\textbf{第二步（证明逆否命题）：}] 假设 $n$ 为奇数，设 $n=2k+1$（$k\in\mathbb{Z}$），则
			\[
				\begin{aligned}
					n^2 &= (2k+1)^2 \\
					&= 4k^2+4k+1 \\
					&= 2(2k^2+2k)+1,
			\end{aligned}
			\]
			即 $n^2$ 为奇数。
			\par\textbf{理由：}
			直接计算证得奇数平方必为奇数，故逆否命题真。

		\item[\textbf{第三步（等价回归）：}] 由于原命题与其逆否命题等价，逆否命题为真，故原命题 $p$ 也为真。
			\par\textbf{理由：}
			原命题 $\equiv$ 逆否命题，所以同真。
	\end{description}

	\textbf{\textcolor{CExample}{关键结论：}}
	\textbf{有些命题直接证明困难时，可转化为证逆否命题，利用等价性达到论证目的。}
\end{examplebox}
